\chapter{Phân tích và thiết kế hệ thống đề xuất}

\section{Thiết kế theo tư duy decision architecture}

Sau khi thiết lập nền phương pháp ở Chương 2, Chương 3 chuyển trọng tâm sang decision architecture: mỗi quyết định thiết kế phải trả lời rõ bốn câu hỏi \textit{(bài toán gì, chọn gì, loại gì, đánh đổi gì)}. Cách tiếp cận này giúp kiến trúc không bị xem như tập hợp công nghệ rời rạc, mà là một chuỗi quyết định kỹ thuật có khả năng kiểm chứng và truy vết.

Trong bối cảnh JobConnect MVP, decision architecture cần giải đồng thời hai mâu thuẫn: (i) muốn tăng chất lượng matching nhưng vẫn giữ chi phí vận hành chấp nhận được, và (ii) muốn tăng tự động hóa nhưng không làm mất khả năng giải trình trong các quyết định tuyển dụng. Do đó, các lựa chọn ở chương này ưu tiên tính module hóa, kiểm soát failure mode, và khả năng rollback theo tầng.

\section{Bối cảnh ràng buộc và yêu cầu thiết kế}

\subsection{Ràng buộc dữ liệu}

Dữ liệu CV/JD là dữ liệu nửa cấu trúc với độ nhiễu cao: parser có thể thiếu trường, mô tả kỹ năng không nhất quán, và nhiều biến thể ngôn ngữ. Ràng buộc dữ liệu này đòi hỏi kiến trúc phải chấp nhận ``không chắc chắn có cấu trúc'', tức là mỗi bước xử lý cần có trạng thái chất lượng và đường thoát khi độ tin cậy giảm.

\subsection{Ràng buộc vận hành}

Hệ thống cần phục vụ các thao tác có tần suất cao (tìm kiếm, so khớp, cập nhật trạng thái), vì vậy độ trễ và độ ổn định quan trọng ngang chất lượng xếp hạng. Một kiến trúc chỉ tối ưu chất lượng offline nhưng thiếu observability và cơ chế suy giảm mềm sẽ khó vận hành thực tế.

\subsection{Ràng buộc tổ chức và governance}

Quyết định tuyển dụng có tác động cao đến ứng viên và nhà tuyển dụng, nên hệ thống phải hỗ trợ giải thích ở mức nghiệp vụ và lưu dấu vết quyết định. Điều này buộc kiến trúc tách biệt rõ ``khuyến nghị'' khỏi ``quyết định cuối cùng'' để tránh tự động hóa quá mức không kiểm soát.

\section{Thiết kế tổng thể theo lớp trách nhiệm}

\begin{figure}[H]
\centering
\fbox{\parbox[c][5.4cm][c]{0.9\textwidth}{\centering
TODO: upload image\\
\textbf{Hình 3.1.} Kiến trúc tổng thể JobConnect MVP theo lớp trách nhiệm
}}
\caption{Placeholder sơ đồ kiến trúc tổng thể để thay bằng hình chuẩn dự án}
\label{fig:jobconnect-architecture-placeholder}
\end{figure}

Kiến trúc được tách thành sáu lớp chính:

\begin{itemize}
  \item \textbf{API Contract Layer}: chuẩn hóa đầu vào/đầu ra và lỗi nghiệp vụ.
  \item \textbf{Identity \& Access Layer}: xác thực, phân vai, giới hạn hành động theo trạng thái.
  \item \textbf{Document Processing Layer}: ingest, parse, normalize, review.
  \item \textbf{Matching Layer}: candidate generation, hard filter, scoring, reranking, explanation.
  \item \textbf{ATS Workflow Layer}: ứng tuyển, lời mời, chuyển trạng thái, notification hooks.
  \item \textbf{Observability \& Audit Layer}: logging có cấu trúc, metric, audit trail, cảnh báo.
\end{itemize}

Cách tách lớp này tạo lợi ích quan trọng: thay đổi trong matching logic không bắt buộc sửa workflow semantics; thay đổi UI/consumer không buộc rewrite pipeline xử lý tài liệu.

\section{Design Decision Log (DDL)}

\begin{table}[H]
\caption{Design Decision Log cho kiến trúc JobConnect MVP}
\centering
\begin{tabularx}{\textwidth}{p{0.07\textwidth}p{0.28\textwidth}p{0.25\textwidth}X}
\toprule
\textbf{ID} & \textbf{Bài toán} & \textbf{Lựa chọn} & \textbf{Trade-off chính}\
\\
\midrule
D1 & Chuẩn giao tiếp giữa module & Contract-first API & Tăng công thiết kế schema ban đầu, đổi lại giảm coupling dài hạn.\\
D2 & Đưa dữ liệu parser vào matching khi nào & Parse-review-publish 3 trạng thái & Tăng một bước vận hành, đổi lại giảm lỗi lan truyền.\\
D3 & Cách lọc ứng viên & Hard filter trước semantic score & Có thể loại sớm vài trường hợp biên, đổi lại tăng tính nhất quán nghiệp vụ.\\
D4 & Cách sinh candidate & Hybrid lexical+dense & Tăng phức tạp chỉ mục, đổi lại tăng độ phủ ngữ nghĩa.\\
D5 & Cách giải thích kết quả & Explanation theo thành phần & Ít chi tiết mô hình sâu, đổi lại dễ hiểu với người nghiệp vụ.\\
D6 & Ranh giới matching/workflow & Tách khuyến nghị khỏi quyết định & Giảm tự động hóa toàn phần, tăng kiểm soát và auditability.\\
D7 & Xử lý lỗi pipeline & Retry theo bước + dead-letter state & Tăng số trạng thái hệ thống, đổi lại tăng khả năng phục hồi.\\
D8 & Quan sát vận hành & Structured logs + metric theo tầng & Tăng chi phí telemetry, đổi lại rút ngắn thời gian điều tra lỗi.\\
D9 & Nâng cấp thuật toán theo thời gian & Progressive enhancement theo lớp & Tốc độ đạt SOTA chậm hơn, đổi lại giảm rủi ro big-bang refactor.\\
D10 & Governance quyết định tuyển dụng & Human-in-the-loop bắt buộc & Giảm mức tự động, đổi lại phù hợp yêu cầu trách nhiệm giải trình.\\
\bottomrule
\end{tabularx}
\end{table}

\section{Thiết kế chi tiết pipeline dữ liệu CV--JD}

\begin{figure}[H]
\centering
\fbox{\parbox[c][5.2cm][c]{0.9\textwidth}{\centering
TODO: upload image\\
\textbf{Hình 3.2.} Pipeline parse/normalize: ingest $\rightarrow$ parse $\rightarrow$ review $\rightarrow$ publish
}}
\caption{Placeholder pipeline chuẩn hóa dữ liệu phục vụ matching}
\label{fig:pipeline-parse-normalize-placeholder}
\end{figure}

\subsection{Ingest và metadata capture}

Bước ingest ghi nhận tài liệu thô, loại tài liệu, nguồn nộp, thời điểm và người thực hiện. Mục tiêu là tạo ngữ cảnh truy vết tối thiểu để mọi lỗi downstream có thể quay ngược về nguồn.

\subsection{Parse và chuẩn hóa trường}

Bước parse tạo biểu diễn bán cấu trúc (skills, experience spans, education, role signals). Chuẩn hóa trường áp dụng mapping đồng nghĩa và kiểm tra nhất quán nội bộ (ví dụ năm kinh nghiệm âm, trường bắt buộc trống, định dạng thời gian không hợp lệ).

\subsection{Review gate trước publish}

Review gate là điểm kiểm soát chất lượng bắt buộc ở MVP. Nếu độ tin cậy dưới ngưỡng, hồ sơ được giữ ở trạng thái chờ hiệu chỉnh thay vì đưa thẳng vào matching. Quyết định này giảm sai lệch khó phát hiện ở downstream và bảo toàn niềm tin người dùng.

\section{Thiết kế matching hai tầng và hợp nhất điểm}

\begin{figure}[H]
\centering
\fbox{\parbox[c][5.2cm][c]{0.9\textwidth}{\centering
TODO: upload image\\
\textbf{Hình 3.3.} Luồng matching 2 tầng: điều kiện bắt buộc $\rightarrow$ chấm điểm ngữ nghĩa
}}
\caption{Placeholder sơ đồ matching hai tầng của JobConnect}
\label{fig:two-stage-matching-placeholder}
\end{figure}

\subsection{Tầng hard filter}

Hard filter kiểm tra các ràng buộc bất khả thương lượng (eligibility constraints). Ở mức hệ thống, hard filter tạo ``vành đai an toàn'' để semantic model không kéo các cặp không khả thi vào top-$k$.

\subsection{Tầng semantic scoring}

Sau hard filter, semantic scoring xử lý tương đồng ngữ nghĩa đa trường. Thay vì dùng một điểm ``opaque'', JobConnect lưu điểm thành phần để phục vụ explanation và calibration theo miền vị trí.

\subsection{Weighted fusion và calibration}

Điểm hợp nhất được tính bằng trọng số động theo cấu hình vị trí hoặc domain. Cơ chế calibration được đặt tách rời khỏi core scorer để tránh hard-code chính sách chấm điểm vào mô hình biểu diễn.

\section{Pseudo-code cho hard-filter + scoring + explanation}

\begin{lstlisting}[language={},caption={Pseudo-code lõi quyết định matching có giải thích},label={lst:core-matching-decision}]
INPUT: anchor a, candidates C, constraints K, weights W
OUTPUT: ranked list with explanation payload

function CORE_MATCH(a, C, K, W):
    valid <- []
    for c in C:
        hf <- evaluate_hard_filters(a, c, K)
        if hf.pass == false:
            continue

        s_lex <- score_lexical(a, c)
        s_dense <- score_dense(a, c)
        s_rule <- score_rule_adjustment(a, c)

        s_final <- W.lex*s_lex + W.dense*s_dense + W.rule*s_rule

        exp <- {
            "hard_filters": hf.details,
            "component_scores": {
                "lexical": s_lex,
                "dense": s_dense,
                "rule_adjustment": s_rule
            },
            "summary": summarize_reasoning(hf, s_lex, s_dense, s_rule),
            "risk_flags": detect_data_risk(a, c)
        }

        valid.append((c, s_final, exp))

    return sort_desc(valid, by="s_final")
\end{lstlisting}

Pseudo-code này minh họa nguyên tắc ``explanation-by-construction'': giải thích được sinh cùng lúc với điểm, không phải vá bổ sung sau khi xếp hạng xong.

\section{Thiết kế explainability và risk flags}

\begin{figure}[H]
\centering
\fbox{\parbox[c][4.8cm][c]{0.9\textwidth}{\centering
TODO: upload image\\
\textbf{Hình 3.4.} Explainability breakdown: hard constraints, component scores, risk flags
}}
\caption{Placeholder sơ đồ bóc tách giải thích kết quả}
\label{fig:explainability-breakdown-placeholder}
\end{figure}

Explainability payload gồm bốn phần:

\begin{itemize}
  \item \textbf{Hard-filter verdict}: điều kiện nào đạt/không đạt.
  \item \textbf{Component scores}: đóng góp từng thành phần điểm.
  \item \textbf{Narrative summary}: diễn giải ngắn cho người nghiệp vụ.
  \item \textbf{Risk flags}: cảnh báo dữ liệu thiếu, parser uncertainty, trường mâu thuẫn.
\end{itemize}

Cấu trúc này phục vụ hai đối tượng khác nhau: recruiter cần lý do hành động; đội kỹ thuật cần tín hiệu chẩn đoán chất lượng dữ liệu và scorer behavior.

\section{Thiết kế ATS state machine và orchestration}

\begin{figure}[H]
\centering
\fbox{\parbox[c][5.0cm][c]{0.9\textwidth}{\centering
TODO: upload image\\
\textbf{Hình 3.5.} State machine ứng tuyển/invite và các chuyển trạng thái chính
}}
\caption{Placeholder state machine cho ATS-lite workflow}
\label{fig:state-machine-placeholder}
\end{figure}

State machine được thiết kế theo nguyên tắc ``explicit transition contract'': mỗi chuyển trạng thái đều cần actor hợp lệ, precondition rõ, và side effect được ghi audit. Cách tiếp cận này giảm lỗi ``nhảy trạng thái'' và tạo cơ sở phân tích chất lượng quy trình tuyển dụng.

\section{Fault isolation và failure modes}

Để vận hành bền vững, hệ thống cần tách failure mode theo lớp thay vì ``một trạng thái lỗi chung''. Bốn failure mode chính:

\begin{itemize}
  \item \textbf{FM1 -- Parse degradation}: dữ liệu văn bản khó trích xuất hoặc mất trường quan trọng.
  \item \textbf{FM2 -- Retrieval drift}: truy vấn mới lệch phân bố khiến candidate generation suy giảm.
  \item \textbf{FM3 -- Scoring instability}: điểm thành phần dao động bất thường giữa các lần chạy gần nhau.
  \item \textbf{FM4 -- Workflow inconsistency}: chuyển trạng thái nghiệp vụ vi phạm precondition.
\end{itemize}

Mỗi failure mode được gắn cơ chế phản ứng riêng: retry, fallback, quarantine state, hoặc human review escalation. Thiết kế này giúp giảm MTTR và tránh lan lỗi giữa các lớp.

\section{Contract-first API, observability và audit trail}

Contract-first được dùng như lớp cách ly thay đổi: khi chỉnh scorer/explanation internals, API contract vẫn giữ ổn định để consumer không bị gãy luồng. Đây là điều kiện quan trọng để phát triển song song backend và lớp tiêu thụ dữ liệu \cite{hld-api,lld-api}.

Ở lớp observability, hệ thống ghi metric theo tầng (ingest, parse, retrieval, scoring, workflow) và log cấu trúc theo request correlation id. Audit trail lưu actor, thời điểm, hành động và trạng thái trước/sau. Kết hợp ba lớp này cho phép vừa giám sát vận hành, vừa phục vụ kiểm tra hậu kiểm.

\section{Phân tích lựa chọn thay thế và lý do loại trừ}

\subsection{Alternative A: monolithic end-to-end model}

\textbf{Ưu điểm}: đơn giản bề mặt kiến trúc, có thể tăng chất lượng top-$k$ nếu huấn luyện tốt.\
\textbf{Lý do loại trừ}: khó giải thích, khó rollback từng phần, khó chẩn đoán lỗi theo lớp.

\subsection{Alternative B: rule-based heavy system}

\textbf{Ưu điểm}: dễ kiểm soát và dễ audit.\
\textbf{Lý do loại trừ}: kém khả năng tổng quát ngữ nghĩa, chi phí bảo trì rule tăng nhanh theo domain.

\subsection{Alternative C: KG-heavy + multi-agent from day one}

\textbf{Ưu điểm}: tiềm năng suy luận giàu ngữ cảnh và mở rộng tri thức.\
\textbf{Lý do loại trừ}: độ phức tạp mô hình hóa và orchestration quá cao so với vòng lặp MVP hiện tại.

\section{Phân rã failure mode theo lớp và chiến lược phục hồi}

Để kiến trúc có thể vận hành liên tục, mỗi lớp cần một chiến lược phản ứng lỗi độc lập thay vì phụ thuộc một cơ chế chung. Bảng dưới đây mô tả các failure mode chính và cơ chế phản hồi tương ứng.

\begin{table}[H]
\caption{Failure mode và chiến lược phục hồi theo lớp kiến trúc}
\centering
\begin{tabularx}{\textwidth}{p{0.16\textwidth}p{0.26\textwidth}p{0.24\textwidth}X}
\toprule
\textbf{Lớp} & \textbf{Failure mode} & \textbf{Phản ứng tức thời} & \textbf{Phản ứng dài hạn}\\
\midrule
Document Processing & Parser timeout, OCR lỗi, thiếu trường bắt buộc & Retry có giới hạn, chuyển trạng thái review-required & Cải thiện template parser, thêm kiểm tra tiền xử lý đầu vào.\\
Retrieval & Candidate set rỗng bất thường, index lệch phiên bản & Fallback sang lexical-only hoặc dense-light & Chuẩn hóa quy trình phát hành index, thêm drift alert theo truy vấn.\\
Hard Filter & Logic constraint sai phiên bản policy & Block transition và ghi audit cảnh báo & Thiết kế versioned business rules và regression suite cho constraint.\\
Scoring/Rerank & Điểm dao động mạnh giữa lần chạy gần nhau & Khoá rerank depth, hạ về cấu hình ổn định hơn & Hiệu chỉnh normalization và tách tuning theo domain.\\
Explanation & Payload thiếu thành phần hoặc narrative sai ngữ cảnh & Trả explanation tối thiểu + risk flag & Đồng bộ schema explanation với pipeline scoring bằng contract test.\\
Workflow & Nhảy trạng thái không hợp lệ, race condition thao tác & Từ chối chuyển trạng thái, yêu cầu retry có điều kiện & Thiết kế optimistic lock/version check cho entity trạng thái.\\
Observability & Thiếu log correlation hoặc metric gián đoạn & Gắn correlation id bắt buộc ở gateway & Chuẩn hóa telemetry contract và health checks cho pipeline quan sát.\\
\bottomrule
\end{tabularx}
\end{table}

Cách phân rã này giúp cô lập lỗi theo tầng, rút ngắn thời gian điều tra và giảm xác suất ``chữa nhầm lớp''. Trong thực tế, khả năng khoanh vùng lỗi nhanh thường quyết định độ tin cậy hệ thống không kém chất lượng mô hình.

\section{Khung testability cho decision architecture}

Một kiến trúc đúng chỉ có giá trị khi kiểm thử được. Với JobConnect, testability được thiết kế theo ba trục:

\begin{itemize}
  \item \textbf{Contract test}: kiểm tra tính ổn định API và explanation schema khi thay scorer.
  \item \textbf{Scenario test}: kiểm tra chuỗi nghiệp vụ từ ingest đến workflow action.
  \item \textbf{Resilience test}: kiểm tra hành vi hệ thống khi từng lớp suy giảm cục bộ.
\end{itemize}

\begin{lstlisting}[language={},caption={Pseudo-code kiểm thử resilience theo lớp},label={lst:resilience-test-harness}]
function RUN_RESILIENCE_TEST(test_case):
    setup_seed_data(test_case.data_profile)
    inject_fault(test_case.layer, test_case.fault_type)

    result <- execute_end_to_end_flow(test_case.flow)

    assert result.system_alive == true
    assert result.fallback_triggered in test_case.allowed_fallbacks
    assert result.audit_trail_complete == true
    assert result.state_consistency == true

    collect_metrics(result)
    rollback_fault_injection()
\end{lstlisting}

Khung này bảo đảm các quyết định ở Design Decision Log không chỉ tồn tại trên tài liệu mà có thể kiểm chứng bằng bằng chứng chạy thực tế.

\section{Bổ sung thiết kế dữ liệu quan sát (observability data model)}

Ngoài dữ liệu nghiệp vụ chính, hệ thống cần một lớp dữ liệu quan sát có cấu trúc. Mỗi sự kiện quan trọng được ghi với tối thiểu các trường: \texttt{event\_time}, \texttt{actor}, \texttt{entity\_type}, \texttt{entity\_id}, \texttt{stage}, \texttt{action}, \texttt{status\_before}, \texttt{status\_after}, \texttt{error\_code}, \texttt{model\_version}, \texttt{index\_version}, \texttt{trace\_id}.

Thiết kế này tạo ba năng lực:

\begin{itemize}
  \item Truy vết toàn tuyến cho một quyết định tuyển dụng.
  \item So sánh hành vi trước/sau khi nâng cấp mô hình hoặc chỉ mục.
  \item Phân tích nguyên nhân gốc khi xuất hiện drift hoặc khiếu nại kết quả.
\end{itemize}

Nếu thiếu observability data model, hệ thống dễ rơi vào trạng thái ``có log nhưng không có bằng chứng'', khiến việc hậu kiểm phụ thuộc suy đoán thay vì dữ liệu.

\section{Kiến trúc mở rộng theo từng pha mà không phá vỡ lõi}

Decision architecture của JobConnect được thiết kế theo hướng mở rộng dần:

\begin{itemize}
  \item \textbf{Pha 1}: củng cố pipeline hiện tại, hoàn thiện telemetry, chuẩn hóa explanation.
  \item \textbf{Pha 2}: thêm reranking sâu có kiểm soát và calibration theo nhóm vị trí.
  \item \textbf{Pha 3}: thử nghiệm KG augmentation hoặc attribution nâng cao trên domain hẹp.
\end{itemize}

\begin{figure}[H]
\centering
\fbox{\parbox[c][4.8cm][c]{0.9\textwidth}{\centering
TODO: upload image\\
\textbf{Hình 3.7.} Lộ trình mở rộng kiến trúc theo pha (không phá vỡ core contract)
}}
\caption{Placeholder lộ trình mở rộng kiến trúc theo pha}
\label{fig:architecture-phased-evolution-placeholder}
\end{figure}

Lợi ích của lộ trình này là giảm rủi ro thay đổi đồng loạt. Mỗi pha chỉ tăng độ phức tạp ở phạm vi có thể quan sát và rollback, giữ ổn định cho lớp workflow và API đã chạy sản phẩm.

\section{Chiến lược versioning policy--model--index và rollback an toàn}

Một rủi ro lớn của hệ thống matching thực chiến là bất đồng bộ phiên bản giữa policy nghiệp vụ, mô hình embedding và chỉ mục truy hồi. Khi ba thành phần thay đổi không đồng bộ, cùng một truy vấn có thể sinh quyết định khác nhau mà không truy vết được nguyên nhân.

Vì vậy, JobConnect cần chiến lược ``triple-version contract'':
\begin{itemize}
  \item \textbf{Policy version}: phiên bản luật lọc cứng, workflow constraint, và ngưỡng quyết định.
  \item \textbf{Model version}: phiên bản embedding/scoring model được áp dụng.
  \item \textbf{Index version}: phiên bản chỉ mục lexical/dense được publish.
\end{itemize}

Mỗi kết quả matching phải mang bộ ba phiên bản này trong metadata. Khi có khiếu nại hoặc regression, hệ thống có thể tái chạy với đúng bộ phiên bản để xác định sai khác nằm ở đâu.

\begin{table}[H]
\caption{Luật kích hoạt rollback theo tín hiệu suy giảm}
\centering
\begin{tabularx}{\textwidth}{p{0.26\textwidth}p{0.18\textwidth}X}
\toprule
\textbf{Tín hiệu quan sát} & \textbf{Rollback scope} & \textbf{Hành động}\\
\midrule
Latency P95 vượt ngưỡng nhiều chu kỳ & Index/serving config & Giảm rerank depth, tạm khóa cấu hình nặng, quay về profile ổn định.\\
Tỷ lệ vi phạm hard constraints tăng & Policy version & Rollback policy về phiên bản trước và mở audit so sánh rule diff.\\
Độ ổn định ranking suy giảm mạnh sau deploy & Model version & Rollback model embedding và đánh dấu lần chạy cần tái hiệu chỉnh calibration.\\
Explanation thiếu thành phần bắt buộc & API contract + scorer & Khóa publish phiên bản mới, trả response fallback có risk flag, sửa scorer pipeline.\\
\bottomrule
\end{tabularx}
\end{table}

Thiết kế này làm tăng chi phí ghi nhận metadata ban đầu, nhưng giảm mạnh thời gian xử lý sự cố và tăng khả năng kiểm chứng quyết định khi cần audit.

\section{Anti-pattern kiến trúc cần tránh trong hệ matching tuyển dụng}

Ngoài các phương án thay thế đã phân tích, cần nêu rõ các anti-pattern triển khai thường gây thất bại:

\begin{itemize}
  \item \textbf{Over-coupled pipeline}: retrieval, filter, scoring và explanation nằm chung một khối logic khó kiểm thử.
  \item \textbf{Silent fallback}: hệ thống fallback nhưng không đánh dấu trong response/audit, gây hiểu nhầm chất lượng.
  \item \textbf{Policy hidden in code}: luật nghiệp vụ nằm rải rác trong hàm scoring, khó cập nhật khi quy trình tuyển dụng đổi.
  \item \textbf{Metric myopia}: tối ưu một metric ranking đơn lẻ mà bỏ qua latency, stability và thao tác vận hành.
\end{itemize}

Các anti-pattern này đặc biệt nguy hiểm trong ATS-lite vì người dùng kỳ vọng quyết định có thể giải thích được ở cấp thao tác. Một pipeline chỉ ``đúng thuật toán'' nhưng ``không đúng vận hành'' vẫn thất bại khi đưa vào quy trình tuyển dụng thật.

\section{Khung ownership và trách nhiệm vận hành theo lớp}

Để tránh vùng xám trách nhiệm, decision architecture cần gắn với ownership rõ theo lớp:

\begin{table}[H]
\caption{Phân tách ownership cho lớp kỹ thuật và lớp nghiệp vụ}
\centering
\begin{tabularx}{\textwidth}{p{0.2\textwidth}p{0.22\textwidth}X}
\toprule
\textbf{Lớp} & \textbf{Owner chính} & \textbf{Trách nhiệm}\\
\midrule
Ingestion/Parsing & Data engineering & Đảm bảo chất lượng dữ liệu đầu vào, phát hiện thiếu trường trọng yếu, kiểm soát parser drift.\\
Retrieval/Ranking & ML/IR engineering & Hiệu chỉnh candidate generation, fusion, rerank; theo dõi chất lượng và độ trễ.\\
Workflow/Policy & Product + backend & Quản trị trạng thái nghiệp vụ, ngưỡng quyết định, và tính hợp lệ chuyển trạng thái.\\
Observability/Audit & Platform/ops & Duy trì telemetry contract, cảnh báo suy giảm, hỗ trợ truy vết sự cố và kiểm toán.\\
\bottomrule
\end{tabularx}
\end{table}

Khi ownership rõ, quá trình nâng cấp theo pha có thể thực hiện song song mà không phá vỡ contract tổng thể. Đây cũng là điều kiện thực tế để roadmap Chương 4 khả thi.

\section{Kết luận chương}

Chương 3 đã cụ thể hóa khung lý thuyết thành hệ quyết định kiến trúc có thể kiểm chứng. Điểm cốt lõi của thiết kế là tách lớp trách nhiệm, kiểm soát failure mode, và tạo explainability ngay trong lõi quyết định. Cấu trúc này cho phép JobConnect vận hành ổn định ở MVP, đồng thời để mở đường nâng cấp dần về thuật toán mà không phá vỡ workflow và contract đã có.
